% =============================================================================
% APPENDIX SK: REF-SKILLS: EBF Slash Commands & Skills Registry
% =============================================================================
% Module: BCM2_00_META
% Version: 1.0 (January 2026)
% Status: Final
% Category: REF (Reference)
% Language: English
%
% This appendix documents all available slash commands (skills) for the EBF
% framework in Claude Code. It serves as the authoritative reference for
% operational automation capabilities.
% APPENDIX SK REF-SKILLS: Claude Code Skill Registry
% =============================================================================
% Category: REF (Reference/Hilfsmittel)
% Module: BCM2_00_META
% Version: 1.0 (January 2026)
% Status: Draft
% Language: EN
%
% Documents the Claude Code skill architecture and EBF integration.
% Single Source of Truth: data/skill-registry.yaml
%
% =============================================================================

% =============================================================================
% CROSS-REFERENCE STRUCTURE
% =============================================================================
%
% CHAPTER LINKAGE:
% - Primary Chapter: N/A (Meta-Infrastructure)
% - Secondary Chapters: All chapters using automated workflows
%
% DEPENDENCIES (what this appendix REQUIRES):
% - Appendix DP (REF-DEVPIPE): Developer Pipeline
% - Appendix APL (REF-APPLY): EBF Workflow Reference
% - Appendix BN (REF-SUPERKEY): Database Superkey System
%
% DEPENDENTS (what REQUIRES this appendix):
% - CLAUDE.md (operational reference)
% - All users of Claude Code skills
%
% IMPLEMENTATION:
% - Source: .claude/commands/*.md
% - Index: .claude/commands/SKILLS-INDEX.md
%
% =============================================================================

\chapter{REF-SKILLS: EBF Slash Commands \& Skills Registry}
\label{app:ref-skills}


% -----------------------------------------------------------------------------
% HEADER BLOCK
% -----------------------------------------------------------------------------

\begin{tcolorbox}[colback=blue!5!white, colframe=blue!75!black,
    title=Appendix Metadata]

\textbf{Appendix:} SK REF-SKILLS

\textbf{Category:} REF (Reference)

\textbf{Version:} 1.0 (January 2026)

\textbf{Status:} Final

\textbf{Language:} English

\textbf{Purpose:} Documents all slash commands (skills) available in the EBF Claude Code environment, providing operational reference for automated workflows.

\textbf{Total Skills:} 42

\end{tcolorbox}


% CHAPTER LINKAGE (Required):
% - Primary Chapter: None (Meta-Documentation)
% - Secondary Chapters: All chapters (skills support entire framework)
%
% DEPENDENCIES (what this appendix REQUIRES):
% - Appendix BN (REF-SUPERKEY): Superkey formats (LEAD, PRJ, SKL)
% - Appendix G (Glossary): Term definitions
%
% DEPENDENTS (what REQUIRES this appendix):
% - All skill implementations in .claude/commands/
% - CLAUDE.md Quick Reference
% - SKILLS-INDEX.md
%
% =============================================================================

\chapter{REF-SKILLS: Claude Code Skill Registry}
\label{app:ref-skills}

% -----------------------------------------------------------------------------
% CHAPTER LINKAGE BOX
% -----------------------------------------------------------------------------

\begin{tcolorbox}[colback=purple!5!white, colframe=purple!75!black,
    title=Chapter Linkage]

\textbf{Primary Chapter:} Meta-Infrastructure (no dedicated chapter)

\textbf{Secondary Chapters:}
\begin{itemize}[nosep]
\item All chapters --- skills support EBF workflow execution
\item CLAUDE.md --- operational implementation
\end{itemize}

\textbf{Related Appendices:}
\begin{itemize}[nosep]
\item Appendix DP (REF-DEVPIPE): Developer Pipeline architecture
\item Appendix APL (REF-APPLY): EBF Workflow steps
\item Appendix BN (REF-SUPERKEY): Database identifiers
\textbf{Primary Chapter:} Meta-Documentation (Framework Infrastructure)
\begin{itemize}[nosep]
\item This appendix documents automation infrastructure
\item Skills support all 25 chapters through automated workflows
\end{itemize}

\textbf{Secondary Chapters:}
\begin{itemize}[nosep]
\item Chapter 17--20: Intervention Design --- \texttt{/design-intervention}, \texttt{/intervention-manage}
\item Chapter 5--13: Core Framework --- \texttt{/design-model}, \texttt{/find-model}
\item All Chapters: Documentation --- \texttt{/compile}, \texttt{/new-chapter}, \texttt{/new-appendix}
\end{itemize}

\textbf{Reading Guide:}
\begin{itemize}[nosep]
\item For skill usage: See \texttt{.claude/commands/*.md}
\item For SSOT: See \texttt{data/skill-registry.yaml}
\item For quick reference: See CLAUDE.md Slash Commands section
\end{itemize}

\end{tcolorbox}


% -----------------------------------------------------------------------------
% ABSTRACT
% -----------------------------------------------------------------------------

\begin{abstract}
This appendix provides the authoritative registry of all slash commands (skills)
available in the EBF Claude Code environment. Skills are automated workflows
that execute specific EBF operations---from model design to document compilation
to project management. The registry documents 40 skills across 11 categories:
Data Management, Literature, Document Creation, Quality Assurance, Build \&
Compilation, Behavioral Modeling, Customer Strategy, Case \& Intervention
Management, and Utilities. Each skill entry includes purpose, usage syntax,
parameters, and integration with other EBF components.
\end{abstract}


This appendix documents the Claude Code skill architecture for the Evidence-Based Framework (EBF).
It defines the Single Source of Truth (SSOT) structure in \texttt{skill-registry.yaml},
categorizes 42 skills across 8 functional domains, and specifies EBF integration requirements
including appendix references, chapter linkages, and 10C dimension mappings.
The registry enables automated project workflows while maintaining framework consistency.
\end{abstract}

% -----------------------------------------------------------------------------
% QUICK REFERENCE BOX
% -----------------------------------------------------------------------------

\begin{tcolorbox}[colback=blue!5!white, colframe=blue!75!black,
    title=Quick Reference: Skills Registry]

\textbf{Core Question:} What automated workflows are available in EBF?

\textbf{Total Skills:} 42 slash commands

\textbf{Categories:}
\begin{itemize}[nosep]
\item Data Management (1)
\item Literature Management (1)
\item Document Creation (3)
\item Quality Assurance (2)
\item Build \& Compilation (3)
\item Behavioral Modeling (2)
\item Customer Strategy (5)
\item Case \& Intervention (6)
\item Sales Pipeline (9)
\item Utilities (3)
\item Client Projects (2)
\end{itemize}

\textbf{Implementation:} \texttt{.claude/commands/*.md}

\textbf{Index:} \texttt{.claude/commands/SKILLS-INDEX.md}

\end{tcolorbox}


% -----------------------------------------------------------------------------
% SCOPE BOX
% -----------------------------------------------------------------------------

\begin{tcolorbox}[
    colback=orange!5!white,
    colframe=orange!75!black,
    title={\textbf{Appendix Scope: Ziel / In-Scope / Out-of-Scope / Constraints}},
    fonttitle=\bfseries
]

\textbf{Ziel (Objective):}
\begin{itemize}[nosep]
    \item Provide authoritative documentation of all EBF slash commands
\end{itemize}

\textbf{In-Scope:}
\begin{itemize}[nosep]
    \item Complete skill inventory (42 commands)
    \item Usage syntax and parameters
    \item Category organization
    \item Integration with EBF workflow
    \item Examples for each skill
\end{itemize}

\textbf{Out-of-Scope (delegiert an andere Appendices):}
\begin{itemize}[nosep]
    \item Developer implementation details $\rightarrow$ Appendix DP (REF-DEVPIPE)
    \item EBF workflow theory $\rightarrow$ Appendix APL (REF-APPLY)
    \item Database schemas $\rightarrow$ Appendix BN (REF-SUPERKEY)
\end{itemize}

\textbf{Constraints:}
\begin{itemize}[nosep]
    \item Skills must be implemented in \texttt{.claude/commands/}
    \item New skills require documentation update here
    \item Skill naming follows kebab-case convention
    title=Quick Reference: Skill Architecture]

\textbf{SSOT:} \texttt{data/skill-registry.yaml}

\textbf{Superkey Format:} \texttt{SKL-\{CATEGORY\}-\{NNN\}}

\textbf{8 Categories:}
\begin{itemize}[nosep]
\item \textbf{PRJ} -- Projekt-Setup \& Management (3 skills)
\item \textbf{FRM} -- Framework \& Modellierung (6 skills)
\item \textbf{SAL} -- Sales Pipeline (8 skills)
\item \textbf{CUS} -- Customer Strategy (5 skills)
\item \textbf{DOC} -- Dokumentation \& LaTeX (6 skills)
\item \textbf{VAL} -- Validierung \& Qualität (4 skills)
\item \textbf{DAT} -- Daten \& Analyse (3 skills)
\item \textbf{SPE} -- Spezielle Workflows (2 skills)
\end{itemize}

\textbf{Total:} 42 Skills with EBF Integration

\end{tcolorbox}

% -----------------------------------------------------------------------------
% APPENDIX SCOPE BOX
% -----------------------------------------------------------------------------

\begin{tcolorbox}[colback=green!5!white, colframe=green!75!black,
    title=Appendix Scope]

\textbf{Ziel:} Formale Dokumentation der Claude Code Skill-Architektur und deren Integration mit dem EBF Framework.

\textbf{In-Scope:}
\begin{itemize}[nosep]
\item Skill-Registry Struktur und SSOT-Prinzipien
\item EBF-Integration (Appendices, Chapters, 10C Dimensions)
\item Dependency Rules und Workflow Chains
\item Kategorie-Definitionen und Superkey-Formate
\end{itemize}

\textbf{Out-of-Scope:}
\begin{itemize}[nosep]
\item Detaillierte Skill-Implementierung (siehe \texttt{.claude/commands/*.md})
\item Python-Script-Dokumentation (siehe Appendix PY)
\item GitHub Actions (siehe \texttt{.github/workflows/})
\end{itemize}

\textbf{Lieferobjekte:}
\begin{itemize}[nosep]
\item L1: SSOT Schema Definition
\item L2: EBF Integration Specification
\item L3: Dependency Rules
\item L4: Worked Examples
\end{itemize}

\end{tcolorbox}


%%%%%%%%%%%%%%%%%%%%%%%%%%%%%%%%%%%%%%%%%%%%%%%%%%%%%%%%%%%%%%%%%%%%%%%%%%%%%%%
\section{The Fundamental Question}
\label{sec:skills-fundamental}
%%%%%%%%%%%%%%%%%%%%%%%%%%%%%%%%%%%%%%%%%%%%%%%%%%%%%%%%%%%%%%%%%%%%%%%%%%%%%%%

\begin{tcolorbox}[colback=red!5!white, colframe=red!75!black,
    title=The Fundamental Question]
\textbf{What automated workflows are available in the EBF ecosystem, and how
do users invoke them to accomplish specific tasks?}
\end{tcolorbox}

The EBF framework provides extensive automation through ``skills''---slash
commands that execute predefined workflows. These skills transform complex
multi-step processes into single-command operations, ensuring consistency
and reducing cognitive load.


%%%%%%%%%%%%%%%%%%%%%%%%%%%%%%%%%%%%%%%%%%%%%%%%%%%%%%%%%%%%%%%%%%%%%%%%%%%%%%%
\section{Theory: Skill Architecture}
\label{sec:skills-theory}
%%%%%%%%%%%%%%%%%%%%%%%%%%%%%%%%%%%%%%%%%%%%%%%%%%%%%%%%%%%%%%%%%%%%%%%%%%%%%%%

The skill system follows a \textbf{command-workflow pattern} that maps user intent
to automated execution:

\begin{equation}
\text{Skill}_i: \text{Intent} \times \text{Parameters} \rightarrow \text{Workflow}_i \rightarrow \text{Output}
\end{equation}

\textbf{Design Principles:}
\begin{enumerate}
    \item \textbf{Single Responsibility}: Each skill handles one specific task type
    \item \textbf{Composability}: Skills can be chained in larger workflows
    \item \textbf{Idempotency}: Re-running a skill produces consistent results
    \item \textbf{Transparency}: All skill actions are logged and reversible
\end{enumerate}

\textbf{Skill Categories} follow the EBF workflow phases (Steps 0-9), ensuring
complete coverage of the behavioral economics project lifecycle.


%%%%%%%%%%%%%%%%%%%%%%%%%%%%%%%%%%%%%%%%%%%%%%%%%%%%%%%%%%%%%%%%%%%%%%%%%%%%%%%
\section{Results: Skill Coverage Analysis}
\label{sec:skills-results}
%%%%%%%%%%%%%%%%%%%%%%%%%%%%%%%%%%%%%%%%%%%%%%%%%%%%%%%%%%%%%%%%%%%%%%%%%%%%%%%

\begin{table}[htbp]
\centering
\begin{tabular}{|l|c|l|}
\hline
\textbf{Category} & \textbf{Count} & \textbf{Primary EBF Step} \\
\hline
Data Management & 1 & Step 1 (Context) \\
Literature Management & 1 & Step 2 (Model) \\
Document Creation & 3 & Step 6 (Report) \\
Quality Assurance & 3 & Step 8 (Quality) \\
Build \& Compilation & 3 & Step 9 (Output) \\
Behavioral Modeling & 3 & Steps 2, 5 (Model, Intervention) \\
Customer Strategy & 5 & Steps 1, 4 (Context, Analysis) \\
Case \& Intervention & 6 & Steps 5, 7 (Intervention, Save) \\
Sales Pipeline & 10 & External (CRM) \\
Utilities & 3 & Cross-cutting \\
Client Projects & 2 & Domain-specific \\
\hline
\textbf{Total} & \textbf{40} & All 10 EBF Steps covered \\
\hline
\end{tabular}
\caption{Skill Distribution by Category}
\end{table}


%%%%%%%%%%%%%%%%%%%%%%%%%%%%%%%%%%%%%%%%%%%%%%%%%%%%%%%%%%%%%%%%%%%%%%%%%%%%%%%
\section{Skills Registry}
\label{sec:skills-registry}
%%%%%%%%%%%%%%%%%%%%%%%%%%%%%%%%%%%%%%%%%%%%%%%%%%%%%%%%%%%%%%%%%%%%%%%%%%%%%%%

% -----------------------------------------------------------------------------
\subsection{Category 1: Data Management}
% -----------------------------------------------------------------------------

\begin{table}[htbp]
\centering
\begin{tabular}{|l|l|p{6cm}|}
\hline
\textbf{Skill} & \textbf{File} & \textbf{Purpose} \\
\hline
\texttt{/classify-papers} & classify-papers.md & Classify extracted papers against paper-sources.yaml \\
\hline
\end{tabular}
\caption{Data Management Skills}
\end{table}


% -----------------------------------------------------------------------------
\subsection{Category 2: Literature Management}
% -----------------------------------------------------------------------------

\begin{table}[htbp]
\centering
\begin{tabular}{|l|l|p{6cm}|}
\hline
\textbf{Skill} & \textbf{File} & \textbf{Purpose} \\
\hline
\texttt{/generate-paper} & generate-paper.md & Generate formal papers in Fehr/Thaler/Kahneman style \\
\hline
\end{tabular}
\caption{Literature Management Skills}
\end{table}


% -----------------------------------------------------------------------------
\subsection{Category 3: Document Creation}
% -----------------------------------------------------------------------------

\begin{table}[htbp]
\centering
\begin{tabular}{|l|l|p{6cm}|}
\hline
\textbf{Skill} & \textbf{File} & \textbf{Purpose} \\
\hline
\texttt{/new-chapter} & new-chapter.md & Create new chapter with template compliance \\
\texttt{/new-appendix} & new-appendix.md & Create new appendix with cross-references \\
\texttt{/doc} & document-production.md & Document Production Workflow (DPW) \\
\hline
\end{tabular}
\caption{Document Creation Skills}
\end{table}


% -----------------------------------------------------------------------------
\subsection{Category 4: Quality Assurance}
% -----------------------------------------------------------------------------

\begin{table}[htbp]
\centering
\begin{tabular}{|l|l|p{6cm}|}
\hline
\textbf{Skill} & \textbf{File} & \textbf{Purpose} \\
\hline
\texttt{/check-compliance} & check-compliance.md & Check chapter/appendix template compliance \\
\texttt{/validate} & validate.md & Run all validation checks \\
\texttt{/integration-test} & integration-test.md & 9-test integration suite for framework validation \\
\hline
\end{tabular}
\caption{Quality Assurance Skills}
\end{table}


% -----------------------------------------------------------------------------
\subsection{Category 5: Build \& Compilation}
% -----------------------------------------------------------------------------

\begin{table}[htbp]
\centering
\begin{tabular}{|l|l|p{6cm}|}
\hline
\textbf{Skill} & \textbf{File} & \textbf{Purpose} \\
\hline
\texttt{/compile} & compile.md & Compile LaTeX to PDF with latexmk \\
\texttt{/convert} & convert.md & Convert between LaTeX, Word, Markdown \\
\texttt{/build-all} & build-all.md & Batch compile all PDFs \\
\hline
\end{tabular}
\caption{Build \& Compilation Skills}
\end{table}


% -----------------------------------------------------------------------------
\subsection{Category 6: Behavioral Modeling}
% -----------------------------------------------------------------------------

\begin{table}[htbp]
\centering
\begin{tabular}{|l|l|p{6cm}|}
\hline
\textbf{Skill} & \textbf{File} & \textbf{Purpose} \\
\hline
\texttt{/design-model} & design-model.md & 9-step EEE workflow for 10C models \\
\texttt{/design-intervention} & design-intervention.md & 20-field schema for interventions \\
\texttt{/find-model} & find-model.md & Search model registry \\
\hline
\end{tabular}
\caption{Behavioral Modeling Skills}
\end{table}


% -----------------------------------------------------------------------------
\subsection{Category 7: Customer Strategy}
% -----------------------------------------------------------------------------

\begin{table}[htbp]
\centering
\begin{tabular}{|l|l|p{6cm}|}
\hline
\textbf{Skill} & \textbf{File} & \textbf{Purpose} \\
\hline
\texttt{/new-customer} & new-customer.md & Create customer database in $<$1 min \\
\texttt{/apply-models} & apply-models.md & Run RPM, MCSM, OSM, CAM models \\
\texttt{/sensitivity-analysis} & sensitivity-analysis.md & What-if parameter analysis \\
\texttt{/board-presentation} & board-presentation.md & Generate 10-slide executive deck \\
\texttt{/replicate-customer} & replicate-customer.md & Clone from ALPLA template \\
\hline
\end{tabular}
\caption{Customer Strategy Skills}
\end{table}


% -----------------------------------------------------------------------------
\subsection{Category 8: Case \& Intervention Registry}
% -----------------------------------------------------------------------------

\begin{table}[htbp]
\centering
\begin{tabular}{|l|l|p{6cm}|}
\hline
\textbf{Skill} & \textbf{File} & \textbf{Purpose} \\
\hline
\texttt{/case} & case.md & Query 10C-indexed case database \\
\texttt{/case-manage} & case-manage.md & Find \& create cases \\
\texttt{/intervention} & intervention.md & Query intervention registry \\
\texttt{/intervention-manage} & intervention-manage.md & Create, update, close projects \\
\texttt{/project-setup} & project-setup.md & Full project setup workflow \\
\texttt{/phase6-manage} & phase6-manage.md & Long-term outcome tracking \\
\hline
\end{tabular}
\caption{Case \& Intervention Registry Skills}
\end{table}


% -----------------------------------------------------------------------------
\subsection{Category 9: Sales Pipeline}
% -----------------------------------------------------------------------------

\begin{table}[htbp]
\centering
\begin{tabular}{|l|l|p{6cm}|}
\hline
\textbf{Skill} & \textbf{File} & \textbf{Purpose} \\
\hline
\texttt{/pipeline-summary} & pipeline-summary.md & Compact pipeline overview \\
\texttt{/pipeline-report} & pipeline-report.md & Detailed pipeline report \\
\texttt{/forecast} & forecast.md & Revenue forecast (Best/Expected/Worst) \\
\texttt{/lead-card} & lead-card.md & Display lead details \\
\texttt{/lead-add} & lead-add.md & Add new lead \\
\texttt{/lead-update} & lead-update.md & Update lead status \\
\texttt{/new-lead} & new-lead.md & Quick lead creation \\
\texttt{/win-loss} & win-loss.md & Deal analysis \\
\texttt{/action-list} & action-list.md & Pipeline task management \\
\texttt{/followup} & followup.md & Contact tracking \& follow-up management \\
\hline
\end{tabular}
\caption{Sales Pipeline Skills}
\end{table}


% -----------------------------------------------------------------------------
\subsection{Category 10: Utilities}
% -----------------------------------------------------------------------------

\begin{table}[htbp]
\centering
\begin{tabular}{|l|l|p{6cm}|}
\hline
\textbf{Skill} & \textbf{File} & \textbf{Purpose} \\
\hline
\texttt{/r-score} & r-score.md & LLMMC to R-Score pipeline \\
\texttt{/bayesian-priors} & bayesian-priors.md & Generate Bayesian priors \\
\texttt{/simulate-stakeholder} & simulate-stakeholder.md & 7-persona focus group simulation \\
\hline
\end{tabular}
\caption{Utility Skills}
\end{table}


% -----------------------------------------------------------------------------
\subsection{Category 11: Client Projects}
% -----------------------------------------------------------------------------

\begin{table}[htbp]
\centering
\begin{tabular}{|l|l|p{6cm}|}
\hline
\textbf{Skill} & \textbf{File} & \textbf{Purpose} \\
\hline
\texttt{/innosuisse} & innosuisse.md & BEATRIX project workflow (AUTO-TRIGGER) \\
\texttt{/bfe-project} & bfe-project.md & BFE intervention projects \\
\hline
\end{tabular}
\caption{Client Project Skills}
\end{table}


%%%%%%%%%%%%%%%%%%%%%%%%%%%%%%%%%%%%%%%%%%%%%%%%%%%%%%%%%%%%%%%%%%%%%%%%%%%%%%%
\section{Skill Invocation}
\label{sec:skills-invocation}
%%%%%%%%%%%%%%%%%%%%%%%%%%%%%%%%%%%%%%%%%%%%%%%%%%%%%%%%%%%%%%%%%%%%%%%%%%%%%%%

\subsection{Basic Syntax}

Skills are invoked with a forward slash followed by the skill name:

\begin{verbatim}
/skill-name [arguments]
\end{verbatim}

\subsection{Examples}

\begin{verbatim}
/design-model --mode schnell
/new-customer "CompanyX" 1500 "Europe,APAC"
/compile outputs/paper.tex
/case --domain health --stage action
/pipeline-summary --period week
\end{verbatim}

\subsection{Help}

To see all available skills:
\begin{verbatim}
/help
ls .claude/commands/
\end{verbatim}


%%%%%%%%%%%%%%%%%%%%%%%%%%%%%%%%%%%%%%%%%%%%%%%%%%%%%%%%%%%%%%%%%%%%%%%%%%%%%%%
\section{Integration with EBF Workflow}
\label{sec:skills-integration}
%%%%%%%%%%%%%%%%%%%%%%%%%%%%%%%%%%%%%%%%%%%%%%%%%%%%%%%%%%%%%%%%%%%%%%%%%%%%%%%

Skills integrate with the 10-step EBF workflow as follows:

\begin{table}[htbp]
\centering
\begin{tabular}{|l|l|l|}
\hline
\textbf{EBF Step} & \textbf{Primary Skills} & \textbf{Purpose} \\
\hline
Step 0: Init & \texttt{/project-setup} & Initialize session \\
Step 1: Context & \texttt{/new-customer} & Load context vectors \\
Step 2: Model & \texttt{/design-model}, \texttt{/find-model} & Build/retrieve model \\
Step 3: Parameters & \texttt{/bayesian-priors} & Estimate parameters \\
Step 4: Analysis & \texttt{/apply-models} & Run analysis \\
Step 5: Intervention & \texttt{/design-intervention} & Design interventions \\
Step 6: Report & \texttt{/board-presentation} & Generate outputs \\
Step 7: Save & \texttt{/intervention-manage} & Register in databases \\
Step 8: Quality & \texttt{/validate}, \texttt{/check-compliance} & Quality checks \\
Step 9: Output & \texttt{/compile}, \texttt{/convert} & Final deliverables \\
\hline
\end{tabular}
\caption{Skills mapped to EBF Workflow Steps}
\end{table}


%%%%%%%%%%%%%%%%%%%%%%%%%%%%%%%%%%%%%%%%%%%%%%%%%%%%%%%%%%%%%%%%%%%%%%%%%%%%%%%
\section{Implications}
\label{sec:skills-implications}
%%%%%%%%%%%%%%%%%%%%%%%%%%%%%%%%%%%%%%%%%%%%%%%%%%%%%%%%%%%%%%%%%%%%%%%%%%%%%%%

\textbf{For Practitioners:}
\begin{itemize}
    \item Skills reduce project setup time from hours to minutes
    \item Standardized workflows ensure consistent output quality
    \item The 10-step EBF workflow is fully automated end-to-end
\end{itemize}

\textbf{For Framework Development:}
\begin{itemize}
    \item New skills must follow the command-workflow pattern
    \item Skills should target a specific EBF step (Steps 0-9)
    \item Documentation in \texttt{.claude/commands/} is authoritative
\end{itemize}


%%%%%%%%%%%%%%%%%%%%%%%%%%%%%%%%%%%%%%%%%%%%%%%%%%%%%%%%%%%%%%%%%%%%%%%%%%%%%%%
\section{Open Issues}
\label{sec:skills-open-issues}
%%%%%%%%%%%%%%%%%%%%%%%%%%%%%%%%%%%%%%%%%%%%%%%%%%%%%%%%%%%%%%%%%%%%%%%%%%%%%%%

\begin{enumerate}
    \item \textbf{Skill Versioning}: No formal versioning for individual skills
    \item \textbf{Dependency Management}: Inter-skill dependencies not formally specified
    \item \textbf{Error Handling}: Inconsistent error reporting across skills
    \item \textbf{Testing}: No automated test suite for skill functionality
\end{enumerate}

These issues are tracked in the project backlog for future resolution.


%%%%%%%%%%%%%%%%%%%%%%%%%%%%%%%%%%%%%%%%%%%%%%%%%%%%%%%%%%%%%%%%%%%%%%%%%%%%%%%
\section{Summary}
\label{sec:skills-summary}
%%%%%%%%%%%%%%%%%%%%%%%%%%%%%%%%%%%%%%%%%%%%%%%%%%%%%%%%%%%%%%%%%%%%%%%%%%%%%%%

\begin{tcolorbox}[colback=green!5!white, colframe=green!75!black,
    title=Summary: REF-SKILLS]

\textbf{Key Points:}
\begin{itemize}[nosep]
\item 40 slash commands across 11 categories
\item Skills automate complex EBF workflows
\item Implementation in \texttt{.claude/commands/*.md}
\item Index at \texttt{.claude/commands/SKILLS-INDEX.md}
\item Skills map to 10-step EBF workflow
\end{itemize}

\textbf{Cross-References:}
\begin{itemize}[nosep]
\item Appendix DP (REF-DEVPIPE): Implementation details
\item Appendix APL (REF-APPLY): Workflow theory
\item Appendix BN (REF-SUPERKEY): Database identifiers
\item Appendix GLS (Glossary): Symbol definitions
\end{itemize}

\end{tcolorbox}


% -----------------------------------------------------------------------------
% GLOSSARY LINK
% -----------------------------------------------------------------------------

\section{Glossary}

For symbol definitions and terminology, see Appendix GLS (Central Glossary).

Key terms used in this appendix:
\begin{itemize}
\item \textbf{Skill}: A slash command that executes a predefined workflow
\item \textbf{10C}: The 10 CORE dimensions of EBF (WHO, WHAT, HOW, WHEN, WHERE, AWARE, READY, STAGE, HIERARCHY, EIT)
\item \textbf{CVA}: Context Vector Architecture
\item \textbf{EEE}: Model Design Workflow (Appendix EEE)
\end{itemize}


% -----------------------------------------------------------------------------
% REFERENCES
% -----------------------------------------------------------------------------

\section{References}

\nocite{bcm_master}

This appendix references the master bibliography. See \texttt{bibliography/bcm\_master.bib} for the complete reference list.

Primary sources for skill documentation:
\begin{itemize}
\item \texttt{.claude/commands/SKILLS-INDEX.md} --- Authoritative skill index
\item \texttt{CLAUDE.md} --- Operational reference
\end{itemize}


% -----------------------------------------------------------------------------
% VERSION HISTORY
% -----------------------------------------------------------------------------

\section{Version History}

\begin{tabular}{|l|l|l|}
\hline
\textbf{Version} & \textbf{Date} & \textbf{Changes} \\
\hline
1.0 & January 2026 & Initial creation with 42 skills \\
\hline
\end{tabular}

\end{chapter}
% =============================================================================
% SECTION 1: THE FUNDAMENTAL QUESTION
% =============================================================================

\section{The Fundamental Question}
\label{sec:sk-fundamental}

\begin{tcolorbox}[colback=yellow!10!white, colframe=yellow!50!black]
\textbf{Central Question:} How can we ensure that automation tools (skills) remain
consistent with the EBF framework while enabling efficient project workflows?
\end{tcolorbox}

The EBF framework comprises 25 chapters, 56 appendices, and extensive data registries.
Manual navigation and consistency checking would be impractical. Claude Code skills
automate common workflows while maintaining framework integrity through:

\begin{enumerate}
\item \textbf{SSOT Principle:} Single source of truth in \texttt{skill-registry.yaml}
\item \textbf{EBF Integration:} Each skill references relevant appendices, chapters, and 10C dimensions
\item \textbf{Dependency Rules:} Enforced prerequisites (e.g., CVA required for project creation)
\item \textbf{Workflow Chains:} Standard sequences linking skills into complete processes
\end{enumerate}

% =============================================================================
% SECTION 2: SKILL REGISTRY ARCHITECTURE
% =============================================================================

\section{Skill Registry Architecture}
\label{sec:sk-architecture}

\subsection{Documentation Hierarchy}

The skill documentation follows a four-level hierarchy:

\begin{verbatim}
┌─────────────────────────────────────────────────────────────────────────┐
│  SKILL DOCUMENTATION ARCHITECTURE                                       │
├─────────────────────────────────────────────────────────────────────────┤
│                                                                         │
│  data/skill-registry.yaml        ← SSOT (machine-readable)             │
│       ↓                                                                 │
│  .claude/commands/SKILLS-INDEX.md  ← Human-readable overview           │
│       ↓                                                                 │
│  .claude/commands/*.md           ← Detailed skill documentation        │
│       ↓                                                                 │
│  CLAUDE.md (Slash Commands)      ← Quick reference                     │
│       ↓                                                                 │
│  Appendix SK (this document)     ← Formal EBF documentation            │
│                                                                         │
└─────────────────────────────────────────────────────────────────────────┘
\end{verbatim}

\subsection{SSOT Schema}

Each skill entry in \texttt{skill-registry.yaml} contains:

\begin{verbatim}
- id: "SKL-PRJ-001"
  command: "/project-setup"
  name: "Project Setup Workflow"
  description: "Automated workflow for complete project setup"
  category: "PRJ"
  file: ".claude/commands/project-setup.md"
  version: "1.0.0"
  created: "2026-01-28"

  ebf_integration:
    appendices:
      - code: "BN"
        name: "REF-SUPERKEY"
        relationship: "Defines superkey formats"
    chapters: [17, 18, 19]
    10c_dimensions: ["WHEN", "WHERE"]

  dependencies:
    requires_cva: true
    creates: ["LEAD-{NNN}", "PRJ-{KUNDE}-{NNN}_scope.yaml"]
    uses_skills: ["SKL-SAL-001", "SKL-CUS-001"]

  tags: [workflow, project, setup, lead, scope]
\end{verbatim}

% =============================================================================
% SECTION 2.3: AXIOMS
% =============================================================================

\subsection{Axioms: Skill Registry Principles}
\label{sec:sk-axioms}

The skill registry is governed by the following axioms:

\begin{description}
\item[Axiom SK-1 (SSOT Principle)] There exists exactly one authoritative source for skill definitions:
\texttt{data/skill-registry.yaml}. All other documentation derives from this source.

\item[Axiom SK-2 (EBF Integration)] Every skill $s \in S$ must specify its EBF integration:
\[
\text{ebf}(s) = \langle A_s, C_s, D_s \rangle
\]
where $A_s$ = referenced appendices, $C_s$ = referenced chapters, $D_s$ = 10C dimensions.

\item[Axiom SK-3 (CVA Dependency)] For skills in categories PRJ, FRM, CUS:
\[
\text{requires\_cva}(s) = \text{true} \Rightarrow \text{CVA must exist before execution}
\]

\item[Axiom SK-4 (Workflow Composition)] Skills can be composed into workflow chains:
\[
W = s_1 \rightarrow s_2 \rightarrow \ldots \rightarrow s_n
\]
where $\text{output}(s_i) \subseteq \text{input}(s_{i+1})$.

\item[Axiom SK-5 (Superkey Uniqueness)] Every skill has a unique identifier:
\[
\forall s_i, s_j \in S: s_i \neq s_j \Rightarrow \text{id}(s_i) \neq \text{id}(s_j)
\]
Format: \texttt{SKL-\{CATEGORY\}-\{NNN\}}.
\end{description}

% =============================================================================
% SECTION 2.4: KEY RESULTS
% =============================================================================

\subsection{Key Results and Propositions}
\label{sec:sk-results}

\begin{proposition}[Skill Coverage]
The 42 skills in 8 categories provide complete coverage of EBF workflow requirements.
\end{proposition}

\begin{proposition}[CVA Necessity]
For behavioral modeling skills, CVA provides necessary context parameters ($\Psi$) without which
parameter cascades ($\lambda$, $\beta$, $\gamma$) cannot be computed.
\end{proposition}

\begin{proposition}[Workflow Completeness]
The standard project workflow chain (Lead $\rightarrow$ CVA $\rightarrow$ Scope $\rightarrow$ Model $\rightarrow$ Intervention $\rightarrow$ Tracking)
covers the complete project lifecycle from acquisition to outcome measurement.
\end{proposition}

% =============================================================================
% SECTION 3: SKILL CATEGORIES
% =============================================================================

\section{Skill Categories}
\label{sec:sk-categories}

\subsection{PRJ: Projekt-Setup \& Management}

Skills for project lifecycle management from lead to outcome tracking.

\begin{table}[h]
\centering
\begin{tabular}{|l|l|l|}
\hline
\textbf{ID} & \textbf{Command} & \textbf{Description} \\
\hline
SKL-PRJ-001 & \texttt{/project-setup} & Complete project setup workflow \\
SKL-PRJ-002 & \texttt{/intervention-manage} & Project creation and closing \\
SKL-PRJ-003 & \texttt{/phase6-manage} & Long-term outcome tracking \\
\hline
\end{tabular}
\caption{PRJ Category Skills}
\end{table}

\textbf{EBF Connection:} Chapters 17--20 (Intervention Design), Appendix HHH (METHOD-TOOLKIT)

\subsection{FRM: Framework \& Modellierung}

Skills for 10C model design and behavioral intervention specification.

\begin{table}[h]
\centering
\begin{tabular}{|l|l|l|}
\hline
\textbf{ID} & \textbf{Command} & \textbf{Description} \\
\hline
SKL-FRM-001 & \texttt{/design-model} & 9-step behavioral model design \\
SKL-FRM-002 & \texttt{/design-intervention} & 20-field intervention schema \\
SKL-FRM-003 & \texttt{/find-model} & Model registry search \\
SKL-FRM-004 & \texttt{/case} & Case registry query \\
SKL-FRM-005 & \texttt{/case-manage} & Case management \\
SKL-FRM-006 & \texttt{/intervention} & Intervention registry query \\
\hline
\end{tabular}
\caption{FRM Category Skills}
\end{table}

\textbf{EBF Connection:} Chapters 5--13 (Core Framework), Appendices EEE/GGG/FFF (METHOD-DESIGN)

\subsection{Other Categories}

Additional categories include:
\begin{itemize}
\item \textbf{SAL} (8 skills): Sales pipeline management (\texttt{/pipeline-summary}, \texttt{/forecast})
\item \textbf{CUS} (5 skills): Customer strategy (\texttt{/new-customer}, \texttt{/apply-models})
\item \textbf{DOC} (6 skills): Documentation (\texttt{/compile}, \texttt{/new-chapter})
\item \textbf{VAL} (4 skills): Validation (\texttt{/check-compliance}, \texttt{/validate})
\item \textbf{DAT} (3 skills): Data analysis (\texttt{/r-score}, \texttt{/bayesian-priors})
\item \textbf{SPE} (2 skills): Special workflows (\texttt{/innosuisse}, \texttt{/simulate-stakeholder})
\end{itemize}

% =============================================================================
% SECTION 4: DEPENDENCY RULES
% =============================================================================

\section{Dependency Rules}
\label{sec:sk-dependencies}

\subsection{CVA Requirement}

The Context Vector Architecture (CVA) is required for many skills:

\begin{tcolorbox}[colback=red!5!white, colframe=red!75!black,
    title=CVA Required Skills]

\textbf{Without CVA, the following skills cannot execute:}
\begin{itemize}[nosep]
\item \texttt{/project-setup} (SKL-PRJ-001)
\item \texttt{/intervention-manage} (SKL-PRJ-002)
\item \texttt{/phase6-manage} (SKL-PRJ-003)
\item \texttt{/design-model} (SKL-FRM-001)
\item \texttt{/design-intervention} (SKL-FRM-002)
\item \texttt{/apply-models} (SKL-CUS-002)
\item \texttt{/sensitivity-analysis} (SKL-CUS-003)
\item \texttt{/board-presentation} (SKL-CUS-004)
\end{itemize}

\textbf{Reason:} EBF models require context parameters ($\Psi$ dimensions).
Without CVA, parameter cascades ($\lambda$, $\beta$, $\gamma$) cannot be computed.

\end{tcolorbox}

\subsection{Workflow Chains}

Standard project workflow linking multiple skills:

\begin{verbatim}
/new-lead → /new-customer → /project-setup → /design-model
         → /design-intervention → /intervention-manage
         → /phase6-manage
\end{verbatim}

% =============================================================================
% SECTION 5: WORKED EXAMPLES
% =============================================================================

\section{Worked Examples}
\label{sec:sk-examples}

\subsection{Example 1: New Project Setup}

\textbf{Scenario:} Create a new project for Helsana (Sales Strategy Sounding)

\begin{verbatim}
# Step 1: Check/create lead
/project-setup "Helsana" "Sounding Sales-Strategie"

# System checks:
# ✓ Lead found: LEAD-047
# ✓ CVA found: data/customers/helsana/ (STANDARD)
# → Proceed to scope definition

# Step 2: Define scope (interactive)
# - Goal: Evaluate potential for behavioral sales optimization
# - In-Scope: Sales strategy sounding with GL/Head of Sales
# - Out-of-Scope: Product development, IT implementation
# - Deliverables: Sounding notes, Fit assessment, Proposal (optional)

# Output:
# → LEAD-047 updated
# → PRJ-HELSANA-001_scope.yaml created
# → Git committed
\end{verbatim}

\subsection{Example 2: Model Design with EBF Integration}

\textbf{Scenario:} Design behavioral model for insurance sales optimization

\begin{verbatim}
/design-model --mode geführt

# System loads:
# - EBF appendices: EEE, GGG, FFF, AAA, BBB
# - Chapters: 5, 9, 10, 11, 12, 13
# - 10C dimensions: WHO, WHAT, HOW, WHEN, WHERE, AWARE, READY, STAGE

# 9-Step workflow:
# Step 1: Entry Point → Practice-Driven
# Step 2: Scope → Decision (insurance purchase)
# Step 3: Context (Ψ) → GGG defaults for CH insurance
# Step 4: Variables (C, γ) → GGG defaults
# Step 5: Functional Form → EBF standard
# Step 6: Parameters (Θ) → BBB literature values + LLMMC
# Step 7: Predictions → Testable hypotheses
# Step 8: Registry → Save to model-registry.yaml
# Step 9: Output → LaTeX/Markdown/Python
\end{verbatim}

% =============================================================================
% SECTION 6: SUMMARY
% =============================================================================

\section{Summary}
\label{sec:sk-summary}

\begin{tcolorbox}[colback=gray!10!white, colframe=gray!75!black,
    title=Summary: Skill Registry]

\textbf{Key Points:}
\begin{enumerate}
\item \textbf{SSOT:} \texttt{data/skill-registry.yaml} is the single source of truth for all 42 skills
\item \textbf{Categories:} 8 functional domains (PRJ, FRM, SAL, CUS, DOC, VAL, DAT, SPE)
\item \textbf{EBF Integration:} Each skill references appendices, chapters, and 10C dimensions
\item \textbf{Dependencies:} CVA is required for project and model-related skills
\item \textbf{Workflow Chains:} Standard sequences ensure consistent project execution
\end{enumerate}

\textbf{Maintenance:}
\begin{enumerate}
\item Update \texttt{skill-registry.yaml} first (SSOT)
\item Update \texttt{SKILLS-INDEX.md} for human readability
\item Update \texttt{CLAUDE.md} quick reference
\item Update this appendix for formal documentation
\end{enumerate}

\end{tcolorbox}

% =============================================================================
% SECTION 7: CRITICAL FOUNDATIONS
% =============================================================================

\section{Critical Foundations}
\label{sec:sk-foundations}

The skill registry architecture rests on the following foundations:

\begin{enumerate}
\item \textbf{Automation Consistency:} Skills automate repetitive tasks while maintaining
EBF framework consistency. Without centralized skill management, workflow variations
would lead to inconsistent outputs.

\item \textbf{Context Dependency:} Behavioral economics models require context ($\Psi$).
The CVA requirement ensures that skills operating on models have access to necessary
context parameters.

\item \textbf{Traceability:} The superkey system (SKL-, LEAD-, PRJ-) enables full
traceability from skill execution to data artifacts.

\item \textbf{Composability:} Skills are designed for composition. The standard workflow
chain demonstrates how atomic skills combine into complete project lifecycles.
\end{enumerate}

\subsection{Potential Objections}

\begin{description}
\item[Objection 1:] ``Skills add complexity without value.'' \\
\textbf{Response:} Skills reduce cognitive load by encapsulating multi-step workflows.
The 9-step model design workflow would otherwise require manual navigation of 6+ appendices.

\item[Objection 2:] ``SSOT creates single point of failure.'' \\
\textbf{Response:} Version control (git) provides full history and recovery capability.
The SSOT principle prevents inconsistencies that would arise from distributed definitions.
\end{description}

% =============================================================================
% SECTION 8: OPEN ISSUES
% =============================================================================

\section{Open Issues and Research Agenda}
\label{sec:sk-open-issues}

\begin{enumerate}
\item \textbf{Skill Versioning:} Current skills use simple version numbers (1.0.0).
A more sophisticated versioning scheme may be needed as skills evolve.

\item \textbf{Cross-Skill Dependencies:} The dependency graph is currently implicit.
A formal dependency resolution system could enable parallel skill execution.

\item \textbf{Skill Testing:} No automated testing framework exists for skills.
Integration tests could verify skill behavior against expected outputs.

\item \textbf{Skill Discovery:} Users currently rely on CLAUDE.md or SKILLS-INDEX.md.
An interactive skill discovery system could improve usability.

\item \textbf{Internationalization:} Skills are currently German/English. Full
internationalization may be needed for global deployment.
\end{enumerate}

% =============================================================================
% SECTION 9: GLOSSARY
% =============================================================================

\section{Glossary of Terms}
\label{sec:sk-glossary}

For comprehensive definitions, see Appendix G (Central Glossary).

\begin{description}
\item[CVA] Context Vector Architecture --- structured customer profile (30--700+ factors)
\item[SSOT] Single Source of Truth --- authoritative data source
\item[Skill] Claude Code slash command with defined workflow
\item[Superkey] Unique identifier format (e.g., SKL-PRJ-001, LEAD-047)
\item[10C] Ten Core Questions framework (WHO, WHAT, HOW, WHEN, WHERE, AWARE, READY, STAGE, HIERARCHY, EIT)
\end{description}

% =============================================================================
% SECTION 8: REFERENCES
% =============================================================================

\section{References}
\label{sec:sk-references}

This appendix references the master bibliography. See \texttt{bibliography/bcm\_master.bib}
for complete citations.

\nocite{bcm_master}

% Cross-references to other appendices:
\begin{itemize}
\item Appendix BN (REF-SUPERKEY): Superkey format definitions
\item Appendix G (Glossary): Central glossary
\item Appendix EEE (METHOD-DESIGN): Model design workflow
\item Appendix HHH (METHOD-TOOLKIT): Intervention toolkit
\item Appendix VC (METHOD-VALIDATE): Validation framework
\end{itemize}

% =============================================================================
% CROSS-REFERENCE FOOTER
% =============================================================================

\vfill
\hrule
\vspace{0.5em}
\small
\textbf{Cross-References:}
Appendix BN (Superkeys) $\rightarrow$
Appendix G (Glossary) $\rightarrow$
Appendix EEE (Model Design) $\rightarrow$
Appendix HHH (Intervention Toolkit)

\textbf{Data Files:}
\texttt{data/skill-registry.yaml} (SSOT) |
\texttt{.claude/commands/} (implementations) |
\texttt{CLAUDE.md} (quick reference)

\normalsize
